\chapter{Cosa resta, e con quali limiti}
\label{cap:cosa_resta}

Ripercorrendo i capitoli di questa terza parte, il cammino compiuto assomiglia più a una serie di correzioni successive che a una linea diritta verso una soluzione. Prima l'idea di un percorso a costo energetico minimo, corretta dal confronto con il Qhapaq Ñan e con lo studio sui sentieri gallesi di Llobera e Sluckin, che hanno mostrato i limiti di un simile modello per gli spostamenti umani pianificati. Poi la conferma, tramite il modello ARC di Pontzer, che quello stesso modello resta valido per il movimento animale, seguita però dalla scoperta che le specie di maggiore interesse, stambecco e camoscio, non compiono affatto lunghi spostamenti da modellare come un corridoio, un limite che la letteratura ecologica generale riconosce esplicitamente per le specie a raggio ridotto. Poi il tentativo, altrettanto instabile, di affidarsi alla presenza di acqua e vegetazione, arenato sull'impossibilità di ricostruire questi elementi con la precisione necessaria per un territorio alpino profondamente rimodellato nel corso dei millenni. Infine, un metodo che funziona davvero, quello di Kompatscher e Hrozny Kompatscher, ma costruito su basi che il territorio veneto, per ora, non possiede: oltre trecento siti già trovati sul campo.

Questo capitolo prova a tirare le fila, indicando cosa resta di realmente utilizzabile, e con quale onestà vada presentato.

\section{Non un criterio unico, ma una gerarchia}

Il punto fermo più solido, emerso con chiarezza dal decimo capitolo, non è una singola caratteristica morfologica da cercare, ma una gerarchia di condizioni che devono combinarsi. Il caso di Compatsch, sull'Alpe di Siusi, lo mostra bene: un'area di cinquantacinque chilometri quadrati, di cui trentasei pianeggianti, eppure quasi priva di siti, perché mancava la seconda condizione, quella delle direttrici verso territori di caccia adiacenti. Il Puflatsch, lo Spitzbühel e il Monte Piz, ai margini della stessa Alpe, restano vuoti nonostante l'ampio spazio a disposizione. Il Passo Sella, al contrario, con un'area sfruttabile più piccola ma con tre direttrici verso territori diversi, concentra almeno trentotto testimonianze.

Questo significa che il criterio di ricerca giusto, per il territorio veneto, non può fermarsi alla sola morfologia isolata (una sella, un rilievo panoramico), ma deve considerare anche cosa quel punto rende raggiungibile: se un valico o una sella si aprono su un solo versante utile, come Compatsch, il valore predittivo cala drasticamente rispetto a un punto che collega più direttrici, come il Passo Sella o il Passo Pampeago.

\section{Due categorie di luogo, adattate alla morfologia locale}

Fermo restando questo principio gerarchico, i casi discussi nel decimo capitolo confermano l'esistenza di due categorie morfologiche ricorrenti, già anticipate nei capitoli precedenti.

La prima categoria è il valico o la sella: un punto di passaggio che mette in comunicazione due versanti o due bacini, spesso scelto proprio perché offre, allo stesso tempo, un accesso più agevole al territorio circostante e una posizione di controllo visivo su un'area ampia. L'esempio dell'Alpe di Luson, discusso nel decimo capitolo, mostra come un percorso possa procedere sistematicamente da sella a sella, aggirando i dossi, con le selle usate come luoghi di bivacco e i dossi intermedi sistematicamente privi di ritrovamenti.

La seconda categoria, distinta dalla prima, è il punto in posizione dominante con buona visuale su un versante o una vallecola sottostante: non necessariamente un valico di transito, ma un punto di osservazione, funzionale a intercettare la fauna che si muove o pascola più in basso. I siti sulla Cresta di Siusi, discussi nel decimo capitolo, rientrano in questa seconda categoria: punti di avvistamento in posizione dominante, distinti dai campi base collocati più in basso, sul versante.

Questa seconda categoria diventa particolarmente rilevante in territori dalla morfologia diversa da quella delle grandi catene alpine con selle nette tra bacini distinti. La Lessinia, discussa nel terzo capitolo di questo libro come area di possibile ricerca futura, è un altopiano carsico inciso da valli parallele, non una successione di creste e bacini separati come il territorio studiato da Kompatscher e Hrozny Kompatscher. In un ambiente di questo tipo, il criterio del valico di collegamento tra due bacini è meno applicabile per definizione geografica: ha più senso cercare piccoli rilievi panoramici affacciati su una vallecola, sul modello dei punti di avvistamento della Cresta di Siusi, piuttosto che selle di passaggio tra due versanti che, nella morfologia lessinica, semplicemente non esistono nella stessa forma.

Distinguere queste due categorie non è un dettaglio tecnico: significa riconoscere che un unico criterio di ricerca, applicato indiscriminatamente a tutto il territorio, rischia di funzionare bene in una zona e male in un'altra, semplicemente perché la morfologia locale offre punti strategici di natura diversa.

\section{Il riparo e l'acqua, con le soglie quantitative}

Accanto alla posizione, un secondo elemento ricorre con costanza nei siti di tipo campo base discussi in questo libro, dal Riparo Tagliente al Riparo Cogola, da Villabruna a Dalmeri: la presenza di un riparo naturale o comunque di un terreno asciutto e leggermente rilevato. I siti di avvistamento descritti da Kompatscher e Hrozny Kompatscher sono invece all'aperto, in posizione dominante: la protezione fisica non è un requisito per questa categoria. La distinzione tra le due tipologie è quindi rilevante anche per questo criterio: un campo base richiede riparo; un punto di avvistamento richiede visuale.

Il decimo capitolo permette ora di aggiungere una soglia quantitativa a questo criterio, che nei capitoli precedenti restava qualitativo: i campi base si trovano tipicamente a venti-ottanta metri da una fonte d'acqua, non più vicini, perché il terreno immediatamente a ridosso di un corso d'acqua tende a essere troppo umido per un accampamento prolungato. I punti di avvistamento, al contrario, possono trovarsi molto più lontani dall'acqua, talvolta senza alcuna fonte nelle vicinanze, come mostra il caso dei Laghi di Pennes, dove i cacciatori si sono spostati di settecento metri dal valico ideale pur di trovare acqua, accampandosi comunque sul bordo esterno del terrazzo, non sulla riva stessa.

Questo suggerisce che il criterio di ricerca più solido, per il territorio veneto, non sia la sola posizione strategica, ma l'incrocio tra questa posizione, la vicinanza a un riparo o terreno asciutto, e una fonte d'acqua nell'ordine delle poche decine di metri per un ipotetico campo base, con maggiore tolleranza per un ipotetico punto di avvistamento.

\section{Cosa questo criterio può e non può fare}

È importante essere chiari, a questo punto, su cosa un criterio di questo tipo permette di ottenere, e cosa no. Un modello costruito incrociando morfologia del terreno (selle, rilievi panoramici, direttrici multiple), presenza di ripari naturali e soglie di distanza dall'acqua non predice dove si trova un sito preistorico. Produce, più modestamente, un elenco di luoghi che meritano una verifica diretta sul terreno, perché condividono le caratteristiche morfologiche osservate nei siti già noti in contesti alpini comparabili.

La differenza tra queste due cose non è sottile. Un modello predittivo in senso stretto affermerebbe: qui, con alta probabilità, si trova un sito. Un modello di questo tipo, più onesto, afferma soltanto: qui, più che altrove, vale la pena andare a guardare. La differenza è la stessa che separa un'ipotesi da una scoperta, e nessun calcolo da tavolino, per quanto ben costruito, può da solo colmare quella distanza. Solo il lavoro sul campo, la prospezione sistematica che Kompatscher e Hrozny Kompatscher hanno praticato per decenni nella loro area di studio, può trasformare un candidato in un sito effettivamente confermato.

Va aggiunta un'ulteriore cautela, emersa proprio dai casi discussi nel decimo capitolo: anche un luogo che soddisfa tutti i criteri sulla carta, come le aree ai margini dell'Alpe di Siusi, può risultare vuoto, mentre un luogo apparentemente meno favorevole può concentrare decine di testimonianze, come il Passo Sella. Questo non significa che i criteri siano inutili: significa che vanno applicati come strumento di prioritizzazione (dove guardare prima), non come garanzia di risultato.

\section{Perché questa modestia è, alla fine, un punto di forza}

Presentare i limiti di un simile criterio non è un modo per sminuire il lavoro fatto in questa parte del libro. È, al contrario, la condizione necessaria perché quel lavoro sia davvero utile a chi, in futuro, volesse metterlo alla prova sul territorio. Un modello che promettesse più di quanto può mantenere rischierebbe di orientare male gli sforzi di ricerca, magari concentrandoli su un singolo luogo indicato con falsa sicurezza, invece che su un ventaglio ragionevole di candidati da verificare con pari attenzione, come hanno fatto per decenni Kompatscher e Hrozny Kompatscher nella loro area di studio.

Con questo criterio in mano, per quanto modesto nelle sue pretese, è possibile tornare finalmente alla domanda con cui questo libro si è aperto: perché le Alpi venete sembrano così povere di arte rupestre rispetto alla Región de Atacama, e dove converrebbe guardare per primi, se si volesse davvero mettere alla prova il sospetto di un research bias discusso nella prima parte. È il compito dell'ultima parte di questo libro.
